\documentclass[12pt,a4paper]{article}
\usepackage[utf8]{inputenc}
\usepackage[T1]{fontenc}
\usepackage[italian]{babel}
\usepackage{geometry}
\usepackage{amsmath}
\usepackage{amssymb}
\usepackage{xcolor}
\usepackage{mdframed}
\usepackage{listings}
\usepackage{hyperref}
\usepackage{bookmark}
\usepackage{booktabs}
\usepackage{array}
\usepackage{multirow}
\usepackage{enumitem}

\geometry{top=2cm, bottom=2cm, left=2.5cm, right=2.5cm}

\hypersetup{
    colorlinks=true,
    linkcolor=blue!70!black,
    urlcolor=blue
}

% Box definizione (grigio)
\newmdenv[
    backgroundcolor=gray!10,
    linecolor=blue!60!black,
    linewidth=1.5pt,
    roundcorner=5pt,
    innertopmargin=6pt,
    innerbottommargin=6pt
]{defbox}

% Box formula (teal)
\newmdenv[
    backgroundcolor=teal!5,
    linecolor=teal!70!black,
    linewidth=1.5pt,
    roundcorner=5pt,
    innertopmargin=6pt,
    innerbottommargin=6pt
]{formulabox}

% Box codice (verde scuro)
\newmdenv[
    backgroundcolor=green!5,
    linecolor=green!50!black,
    linewidth=1pt,
    roundcorner=3pt,
    innertopmargin=4pt,
    innerbottommargin=4pt
]{codebox}

% Box risultati (arancione)
\newmdenv[
    backgroundcolor=orange!8,
    linecolor=orange!70!black,
    linewidth=1.5pt,
    roundcorner=5pt,
    innertopmargin=6pt,
    innerbottommargin=6pt
]{resultbox}

% Stile listings Python
\lstdefinestyle{python}{
    language=Python,
    basicstyle=\ttfamily\footnotesize,
    keywordstyle=\color{blue!80!black}\bfseries,
    commentstyle=\color{green!50!black}\itshape,
    stringstyle=\color{red!70!black},
    numbers=left,
    numberstyle=\tiny\color{gray},
    numbersep=8pt,
    breaklines=true,
    showstringspaces=false,
    frame=none,
    tabsize=4,
    backgroundcolor=\color{gray!5}
}

\title{\textbf{TLN -- Parte 3: Di Caro}\\[0.3em]
\large Similarit\`a, Topic Modeling e RAG\\[0.2em]
\normalsize A.A. 2025/26 -- UniTO Magistrale Informatica}
\author{Riassunto per lo studio}
\date{}

\begin{document}
\maketitle
\tableofcontents
\newpage

%% ============================================================
\section{Introduzione: Contenuto vs. Forma Lessicale}
%% ============================================================

\begin{defbox}
\textbf{Due direzioni nel lessico:}
\begin{itemize}
  \item \textbf{Semasiologica} (Forma $\to$ Contenuto): dato un lemma, cerco il suo significato. Direzione classica del dizionario.
  \item \textbf{Onomasiologica} (Contenuto $\to$ Forma): data una descrizione/definizione, trovo la parola corretta. Problema di \emph{content-to-form retrieval}.
\end{itemize}
\end{defbox}

\medskip
Il percorso del corso copre tre macro-temi:
\begin{enumerate}
  \item \textbf{Similarit\`a lessicale/semantica} (Lab 1--2): misure Jaccard e TF-IDF coseno su definizioni umane
  \item \textbf{Topic Modeling neurale} (Lab 4): pipeline BERTopic per scoperta automatica di argomenti
  \item \textbf{RAG ibrido} (Lab 5): retrieval aumentato con FAISS semantico + BM25 lessicale + RRF + LLM
\end{enumerate}

%% ============================================================
\section{I Quattro Tipi Concettuali}
%% ============================================================

Nei lab 1 e 2 si lavora su quattro concetti scelti lungo due assi: \textbf{concretezza} (concreto/astratto) e \textbf{specificit\`a} (generico/specifico).

\begin{center}
\begin{tabular}{lcccc}
\toprule
\textbf{Tipo} & \textbf{Sigla} & \textbf{Esempio} & \textbf{Concretezza} & \textbf{Specificit\`a}\\
\midrule
Concrete Generic   & CG & Albero (Tree)  & Concreto  & Generico \\
Concrete Specific  & CS & Teiera (Teapot) & Concreto  & Specifico \\
Abstract Generic   & AG & Musica (Music)  & Astratto  & Generico \\
Abstract Specific  & AS & Etica (Ethics)  & Astratto  & Specifico \\
\bottomrule
\end{tabular}
\end{center}

\begin{defbox}
\textbf{Caratteristiche per tipo:}
\begin{itemize}
  \item \textbf{CG (Albero)}: referente fisico, classe ampia. Definizioni usano vocabolario condiviso (radici, tronco, rami, foglie). Alta coerenza cross-individuale.
  \item \textbf{CS (Teiera)}: referente fisico specifico, funzione primaria. Vocabolario concentrato su funzione (infusione, versare).
  \item \textbf{AG (Musica)}: nessun referente fisico, classe ampia. Definizioni variano enormemente: operative, sensoriali, emotive.
  \item \textbf{AS (Etica)}: nessun referente fisico, dominio filosofico. Alta influenza culturale e contestuale. Eterogeneit\`a massima.
\end{itemize}
\end{defbox}

%% ============================================================
\section{Similarit\`a Lessicale e Semantica (Lab 1)}
%% ============================================================

\subsection{SimLex -- Indice di Jaccard}

\begin{formulabox}
\[
\text{SimLex}(A, B) = \frac{|A \cap B|}{|A \cup B|}
\]
dove $A$ e $B$ sono insiemi di parole lemmatizzate e filtrate (stopwords rimosse) delle due definizioni.
\end{formulabox}

\textbf{Cosa misura:} sovrapposizione \emph{lessicale} pura -- due testi con stesso significato ma parole diverse $\to$ score basso.

\textbf{Pipeline di preprocessing:}
\begin{enumerate}
  \item Lowercase
  \item Rimozione punteggiatura
  \item Tokenizzazione (NLTK \texttt{word\_tokenize})
  \item Rimozione stopwords (NLTK \texttt{stopwords.words('italian')})
  \item Lemmatizzazione
  \item Conversione in \texttt{set}
\end{enumerate}

\begin{codebox}
\begin{lstlisting}[style=python]
from nltk.corpus import stopwords
from nltk.tokenize import word_tokenize
import re, string

def preprocess_to_set(text, lang='italian'):
    text = text.lower()
    text = re.sub(r'[^\w\s]', '', text)
    tokens = word_tokenize(text, language=lang)
    stop = set(stopwords.words(lang))
    tokens = [t for t in tokens if t not in stop and len(t) > 1]
    return set(tokens)  # set per Jaccard

def simlex(defA, defB):
    A = preprocess_to_set(defA)
    B = preprocess_to_set(defB)
    if not A and not B:
        return 0.0
    return len(A & B) / len(A | B)
\end{lstlisting}
\end{codebox}

\subsection{SimSem -- Coseno TF-IDF}

\begin{formulabox}
\[
\text{SimSem}(A, B) = \cos(\vec{v}_A, \vec{v}_B) = \frac{\vec{v}_A \cdot \vec{v}_B}{||\vec{v}_A|| \cdot ||\vec{v}_B||}
\]
dove $\vec{v}$ \`e la rappresentazione TF-IDF del testo.
\[
\text{TF-IDF}(t, d) = \text{TF}(t,d) \times \log\frac{N}{\text{df}(t)}
\]
\end{formulabox}

\textbf{Cosa misura:} prossimit\`a \emph{semantica} -- cattura sinonimi e termini correlati che il semplice overlap non coglie.

\begin{codebox}
\begin{lstlisting}[style=python]
from sklearn.feature_extraction.text import TfidfVectorizer
from sklearn.metrics.pairwise import cosine_similarity

def simsem(defA, defB):
    corpus = [preprocess_text(defA), preprocess_text(defB)]
    vectorizer = TfidfVectorizer()
    tfidf_matrix = vectorizer.fit_transform(corpus)
    return cosine_similarity(tfidf_matrix[0:1], tfidf_matrix[1:2])[0][0]
\end{lstlisting}
\end{codebox}

\subsection{Confronto tra coppie -- tutte le combinazioni}

Il codice usa \texttt{itertools.combinations} per calcolare similarit\`a tra tutte le coppie di concetti:

\begin{codebox}
\begin{lstlisting}[style=python]
import itertools, csv

concetti = ['Musica', 'Etica', 'Albero', 'Teiera']
definizioni = {c: [...] for c in concetti}  # lista def per c

risultati = []
for c1, c2 in itertools.combinations(concetti, 2):
    for def1, def2 in zip(definizioni[c1], definizioni[c2]):
        sl = simlex(def1, def2)
        ss = simsem(def1, def2)
        risultati.append({'concetto1': c1, 'concetto2': c2,
                          'simlex': sl, 'simsem': ss})

# salva su CSV
with open('risultati_similarita.csv', 'w') as f:
    writer = csv.DictWriter(f, fieldnames=risultati[0].keys())
    writer.writeheader()
    writer.writerows(risultati)
\end{lstlisting}
\end{codebox}

\subsection{Risultati Lab 1}

\begin{resultbox}
\begin{center}
\begin{tabular}{lcc}
\toprule
\textbf{Concetto} & \textbf{SimLex (Jaccard)} & \textbf{SimSem (Coseno)} \\
\midrule
Albero (CG)  & 0.084 & 0.077 \\
Teiera (CS)  & 0.057 & 0.062 \\
Etica (AS)   & 0.034 & 0.041 \\
Musica (AG)  & 0.034 & 0.036 \\
\bottomrule
\end{tabular}
\end{center}
\end{resultbox}

\textbf{Osservazioni chiave:}
\begin{itemize}
  \item Valori molto bassi ovunque (3--8\%): enorme variabilit\`a del linguaggio naturale anche per lo stesso referente.
  \item \textbf{Concreto > Astratto}: i concetti fisici hanno referenti che vincolano il vocabolario.
  \item \textbf{Per i concetti astratti, SimSem > SimLex}: il modello vettoriale allinea sinonimi e termini correlati che la Jaccard pura non vede. Questo dimostra il valore aggiunto della rappresentazione semantica.
  \item Etica: sebbene \texttt{simlex = simsem = 0.034}, le definizioni condividono \emph{struttura concettuale} nonostante parole diverse.
\end{itemize}

%% ============================================================
\section{Content-to-Form: Ricerca Onomasiologica (Lab 2)}
%% ============================================================

\subsection{Il Task}

\begin{defbox}
\textbf{Ricerca onomasiologica}: data una definizione in linguaggio naturale, recuperare il \emph{synset} WordNet corretto.\\[4pt]
Input: \textit{"Branch of philosophy that studies human behavior by trying to define what is right and wrong"} \\
Output: \texttt{ethics.n.01}
\end{defbox}

\`E l'inverso della classica consultazione dizionario: si parte dal concetto/significato per trovare la forma lessicale.

\subsection{Procedura Lab 2}

\begin{enumerate}
  \item \textbf{Selezione candidati}: per ogni concetto target, 5 synset WordNet candidati con relative glosse (incluse forme alternate e iperonimi)
  \item \textbf{Preprocessing}: lowercase, rimozione punteggiatura, filtro stopwords su definizione utente E glosse
  \item \textbf{Vettorizzazione TF-IDF}: l'intero corpus = [definizione\_pulita] + [glosse\_candidate\_pulite]
  \item \textbf{Coseno}: similarit\`a tra definizione e ogni glossa candidata
  \item \textbf{Predizione}: synset con score massimo = risposta del sistema
  \item \textbf{Valutazione}: confronto con synset target
\end{enumerate}

\begin{codebox}
\begin{lstlisting}[style=python]
from sklearn.feature_extraction.text import TfidfVectorizer
from sklearn.metrics.pairwise import cosine_similarity

# Per ogni concetto: definizione umana + 5 glosse candidate
concetti = {
    'Albero': {
        'target': 'tree.n.01',
        'definizioni': [...],  # definizioni degli studenti
        'candidati': {
            'tree.n.01': 'a tall perennial woody plant having...',
            'plant.n.01': 'a living organism lacking the power...',
            # ... altri 3 candidati
        }
    },
    # ... altri concetti
}

def content_to_form(defn, glosse_dict):
    glosse = list(glosse_dict.values())
    synsets = list(glosse_dict.keys())
    corpus = [clean(defn)] + [clean(g) for g in glosse]

    vec = TfidfVectorizer()
    mat = vec.fit_transform(corpus)
    scores = cosine_similarity(mat[0:1], mat[1:])[0]
    best_idx = scores.argmax()
    return synsets[best_idx], scores[best_idx]
\end{lstlisting}
\end{codebox}

\subsection{Struttura Genus-Differentia}

\begin{defbox}
\textbf{Schema Aristotelico}: la migliore strategia definitoria \`e \emph{genus proximum + differentia specifica}:
\begin{itemize}
  \item \textbf{Genus}: categoria sovraordinata (spesso corrisponde all'ipernimo WordNet)
  \item \textbf{Differentia}: propriet\`a che distinguono il concetto dagli altri della stessa classe
\end{itemize}
\textbf{Esempio efficace:} \textit{"Branch of philosophy [genus] that studies moral values [differentia]"}\\
Corrisponde quasi letteralmente alla glossa \texttt{ethics.n.01}: \textit{"the philosophical study of moral values and rules"}
\end{defbox}

\subsection{Risultati Lab 2}

\begin{resultbox}
\begin{center}
\begin{tabular}{lcccc}
\toprule
\textbf{Concetto} & \textbf{Tipo} & \textbf{Accuracy} & \textbf{Score Avg} & \textbf{Score Max} \\
\midrule
Albero (CG)  & Concreto/Generico  & 89.3\% & 0.1025 & 0.2919 \\
Teiera (CS)  & Concreto/Specifico & 71.4\% & 0.1272 & 0.3063 \\
Etica (AS)   & Astratto/Specifico & 42.9\% & 0.1524 & 0.6295 \\
Musica (AG)  & Astratto/Generico  & 32.1\% & 0.1298 & 0.3981 \\
\bottomrule
\end{tabular}
\end{center}

\bigskip
\textbf{Aggregato per dimensione:}
\begin{center}
\begin{tabular}{lcc}
\toprule
\textbf{Dimensione} & \textbf{Accuracy} & \textbf{Score Avg} \\
\midrule
Concreto & 80.4\% & 0.1149 \\
Astratto & 37.5\% & 0.1411 \\
Generico & 60.7\% & 0.1162 \\
Specifico & 57.1\% & 0.1398 \\
\bottomrule
\end{tabular}
\end{center}
\end{resultbox}

\textbf{Analisi:}
\begin{itemize}
  \item Concetti concreti: 2.1$\times$ pi\`u recuperabili degli astratti. Il referente fisico vincola il vocabolario.
  \item Definizione migliore per Etica (score 0.63): ``Branch of philosophy that studies human behavior by trying to define what is right and what is wrong, often influenced by historical and sociocultural context.'' -- ha genus-differentia esplicita.
  \item TF-IDF fallisce su sinonimia: richiede embedding contestuali per i concetti astratti.
\end{itemize}

%% ============================================================
\section{Topic Modeling con BERTopic (Lab 4)}
%% ============================================================

\subsection{Architettura della Pipeline}

\begin{defbox}
\textbf{Pipeline BERTopic} (4 stadi):
\[
\text{Documenti} \xrightarrow{\text{1. Embedding}} \xrightarrow{\text{2. UMAP}} \xrightarrow{\text{3. HDBSCAN}} \xrightarrow{\text{4. c-TF-IDF}} \text{Topic interpretabili}
\]
\end{defbox}

\subsection{Stadio 1: Embedding}

Ogni documento viene rappresentato come vettore denso semantico tramite un \emph{SentenceTransformer}.

\begin{center}
\begin{tabular}{lcc}
\toprule
\textbf{Modello} & \textbf{Dimensioni} & \textbf{Caratteristiche} \\
\midrule
\texttt{thenlper/gte-small} & 384 & Qualit\`a alta, lento su CPU (13.484s) \\
\texttt{all-MiniLM-L6-v2} & 384 & Compatto, veloce su CPU, comparabile \\
\bottomrule
\end{tabular}
\end{center}

\subsection{Stadio 2: UMAP (Riduzione Dimensionale)}

\begin{defbox}
\textbf{UMAP} (Uniform Manifold Approximation and Projection): riduce le dimensioni preservando sia struttura locale che globale.
\begin{itemize}
  \item \textbf{n\_components}: dimensioni output (meno = pi\`u veloce, meno dettaglio semantico)
  \item \textbf{n\_neighbors}: dimensione del vicinato locale (basso = focus locale; alto = focus globale)
  \item \textbf{min\_dist}: distanza minima tra punti (0.0 per clustering; pi\`u alto per visualizzazione)
  \item \textbf{metric}: ``cosine'' per embedding normalizzati
\end{itemize}
Vantaggio su t-SNE: preserva struttura globale oltre a quella locale.
\end{defbox}

\subsection{Stadio 3: HDBSCAN (Clustering)}

\begin{defbox}
\textbf{HDBSCAN} (Hierarchical Density-Based Spatial Clustering of Applications with Noise):
\begin{itemize}
  \item Identifica cluster di forma arbitraria (non solo sferici)
  \item Assegna automaticamente outlier all'etichetta $-1$
  \item \textbf{Non richiede} di specificare il numero di cluster a priori
  \item Parametro chiave: \textbf{min\_cluster\_size} -- dimensione minima cluster
\end{itemize}
\end{defbox}

\subsection{Stadio 4: c-TF-IDF (Estrazione Keywords Topic)}

\begin{formulabox}
\[
c\text{-TF-IDF}(t, c) = \frac{\text{freq}(t, c)}{|c|} \times \log\frac{N_{\text{docs}}}{|\{d : t \in d\}|}
\]
dove $c$ \`e il cluster (topic), non il singolo documento.
\end{formulabox}

Differenza da TF-IDF classico: la frequenza viene calcolata \emph{all'interno del cluster} come contesto, non del singolo documento. Produce keyword interpretabili e specifiche per ogni topic.

\subsection{Configurazioni Testate}

\begin{center}
\begin{tabular}{llccc}
\toprule
 & \textbf{Modello} & \textbf{UMAP} & \textbf{HDBSCAN} & \textbf{Effetto} \\
\midrule
Config 1 (base)   & gte-small   & comp=5, neigh=15  & mcs=50  & Bilanciato \\
Config 2 (fine)   & MiniLM-L6   & comp=10, neigh=10 & mcs=30  & Granulare \\
Config 3 (coarse) & MiniLM-L6   & comp=5, neigh=30  & mcs=150 & Grossolano (fallimento) \\
\bottomrule
\end{tabular}
\end{center}

\subsection{Risultati Comparativi}

\begin{resultbox}
\begin{center}
\begin{tabular}{lccc}
\toprule
\textbf{Metrica} & \textbf{Config 1} & \textbf{Config 2} & \textbf{Config 3} \\
\midrule
\# Topic (escl. $-1$) & 151 & 171 & 49 \\
Outlier count         & 14.357 & 12.747 & 14.899 \\
Outlier \%            & 31.9\% & 28.4\% & 33.1\% \\
Runtime               & $\sim$3h 47m & $\sim$17m 30s & $\sim$11m \\
\midrule
Topic 0 keywords & speech, asr, recognition & dialogue, response, conv. & dialogue, \textbf{the}, \textbf{to}, \textbf{and} \\
\bottomrule
\end{tabular}
\end{center}
\end{resultbox}

\textbf{Analisi failure mode Config 3:}

\begin{itemize}
  \item \texttt{min\_cluster\_size=150}: cluster troppo grandi e tematicamente eterogenei
  \item Il c-TF-IDF non riesce a estrarre segnale discriminativo: le stopwords emergono nel top-5
  \item Solo 49 topic (vs. 151--171 degli altri) con outlier percentuale paradossalmente pi\`u alta
  \item \textbf{Lezione}: parametri troppo grossolani $\to$ perdita di interpretabilit\`a
\end{itemize}

\textbf{Config 2 (raccomandata per CPU):}

\begin{codebox}
\begin{lstlisting}[style=python]
from sentence_transformers import SentenceTransformer
from bertopic import BERTopic
from umap import UMAP
from hdbscan import HDBSCAN

embedding_model = SentenceTransformer("all-MiniLM-L6-v2")
umap_model  = UMAP(n_components=10, n_neighbors=10,
                   min_dist=0.0, metric='cosine')
hdbscan_model = HDBSCAN(min_cluster_size=30,
                        cluster_selection_method='eom',
                        prediction_data=True)

topic_model = BERTopic(
    embedding_model=embedding_model,
    umap_model=umap_model,
    hdbscan_model=hdbscan_model,
    language="english",
    verbose=True
)

topics, probs = topic_model.fit_transform(docs)
print(topic_model.get_topic_info().head(10))
\end{lstlisting}
\end{codebox}

Dataset: \texttt{MaartenGr/arxiv\_nlp} (HuggingFace) -- 44.949 abstract di paper NLP.

%% ============================================================
\section{RAG: Retrieval-Augmented Generation}
%% ============================================================

\subsection{Architettura Generale}

\begin{defbox}
\textbf{RAG} combina un \emph{retriever} (recupera documenti rilevanti) con un \emph{generatore} (LLM che risponde basandosi sul contesto):
\[
\text{Query} \to \underbrace{[\text{Retriever}]}_{\text{top-}k\text{ docs}} \to \underbrace{[\text{Context Assembly}]}_{\text{concatenazione}} \to \underbrace{[\text{LLM}]}_{\text{risposta}} \to \text{Answer}
\]

\textbf{Vantaggi rispetto a LLM standalone:}
\begin{itemize}
  \item Accesso a conoscenza esterna e aggiornabile senza retraining
  \item Riduzione delle \emph{allucinazioni} (risposta ancorata ai documenti)
  \item Possibilit\`a di citare le fonti
  \item Adattabile a domini specifici
\end{itemize}
\end{defbox}

\subsection{Retrieval Semantico con FAISS + SciBERT}

\begin{defbox}
\textbf{FAISS} (Facebook AI Similarity Search): libreria per nearest-neighbor efficiente su vettori densi.
\begin{itemize}
  \item \textbf{IndexFlatIP}: inner product search su vettori normalizzati = cosine similarity
  \item Vettori: SciBERT (\texttt{allenai/scibert\_scivocab\_uncased}), 768 dimensioni
  \item SciBERT: BERT pre-trained su 1.14M paper scientifici (Semantic Scholar)
\end{itemize}
\end{defbox}

\begin{formulabox}
\[
\text{Cosine}(q, d) = \frac{\vec{q} \cdot \vec{d}}{||\vec{q}|| \cdot ||\vec{d}||} \quad \text{(equivalente a inner product su vettori normalizzati)}
\]
\end{formulabox}

\textbf{Punti di forza}: cattura relazioni semantiche e sinonimi\\
\textbf{Limiti}: impreciso su termini tecnici esatti (acronimi, nomi di algoritmi)

\subsection{Retrieval Lessicale con BM25}

\begin{formulabox}
\textbf{BM25} (Best Match 25):
\[
\text{BM25}(D,Q) = \sum_{i} \text{IDF}(q_i) \cdot \frac{f(q_i,D) \cdot (k_1+1)}{f(q_i,D) + k_1 \cdot \left(1 - b + b \cdot \frac{|D|}{\text{avgdl}}\right)}
\]
\[
\text{IDF}(q_i) = \log\frac{N - n(q_i) + 0.5}{n(q_i) + 0.5}
\]
dove $k_1 = 1.5$, $b = 0.75$, $N$ = dimensione corpus, $n(q_i)$ = doc frequency del termine.
\end{formulabox}

\begin{itemize}
  \item \textbf{Saturazione}: l'effetto della frequenza del termine raggiunge un plateau (parametro $k_1$)
  \item \textbf{Normalizzazione per lunghezza}: documenti lunghi non avvantaggiati ingiustamente (parametro $b$)
  \item Molto veloce (nessuna inferenza neurale)
\end{itemize}

\textbf{Punti di forza}: preciso su vocabolario tecnico (acronimi, termini specifici)\\
\textbf{Limiti}: nessuna comprensione semantica (sinonimi = miss), sensibile alla distribuzione del vocabolario

\subsection{Hybrid Search: Reciprocal Rank Fusion (RRF)}

\begin{defbox}
\textbf{Problema}: FAISS e BM25 hanno punteggi non confrontabili (cosine $\in [0,1]$ vs score BM25 $\in \mathbb{R}^+$).\\
\textbf{Soluzione}: combinare le \emph{classifiche} (rankings) anzich\'e i punteggi assoluti.
\end{defbox}

\begin{formulabox}
\[
\text{RRF}(d) = \sum_{r \in R} \frac{1}{k + \text{rank}_r(d)}
\]
dove $R$ = insieme dei retriever, $k = 60$ (costante di smorzamento standard).
\end{formulabox}

\textbf{Perch\'e funziona:}
\begin{itemize}
  \item Robusto all'eterogeneit\`a degli score
  \item Favorisce documenti in cima a pi\`u classifiche (complementariet\`a)
  \item Nessun tuning manuale di pesi relativi
  \item Semplice e interpretabile
\end{itemize}

\textbf{Procedura Lab 5:}
\begin{enumerate}
  \item BM25: top-20 per score lessicale
  \item FAISS: top-20 per cosine similarity semantica
  \item RRF: somma $\frac{1}{k+\text{rank}}$ per ogni documento nelle due liste
  \item Re-ranking per score RRF finale, restituire top-$k$
\end{enumerate}

\begin{codebox}
\begin{lstlisting}[style=python]
from rank_bm25 import BM25Okapi
import faiss, numpy as np

# Inizializzazione
tokenized_docs = [doc.lower().split() for doc in documents]
bm25 = BM25Okapi(tokenized_docs)
faiss_index = faiss.IndexFlatIP(768)
faiss_index.add(embeddings_normalized)

def hybrid_retrieve(query, k=5, k_rrf=60):
    # BM25 retrieval
    bm25_scores = bm25.get_scores(query.lower().split())
    bm25_top = np.argsort(bm25_scores)[::-1][:20]

    # FAISS retrieval
    q_emb = encode_and_normalize(query)
    _, faiss_top = faiss_index.search(q_emb, 20)
    faiss_top = faiss_top[0]

    # Reciprocal Rank Fusion
    rrf_scores = {}
    for rank, idx in enumerate(bm25_top):
        rrf_scores[idx] = rrf_scores.get(idx, 0) + 1/(k_rrf + rank + 1)
    for rank, idx in enumerate(faiss_top):
        rrf_scores[idx] = rrf_scores.get(idx, 0) + 1/(k_rrf + rank + 1)

    sorted_docs = sorted(rrf_scores.items(), key=lambda x: x[1], reverse=True)
    return [documents[i] for i, _ in sorted_docs[:k]]
\end{lstlisting}
\end{codebox}

\subsection{Componenti del Sistema (Lab 5)}

\begin{center}
\begin{tabular}{ll}
\toprule
\textbf{Componente} & \textbf{Implementazione} \\
\midrule
Dataset & \texttt{MaartenGr/arxiv\_nlp} (44.949 abstract NLP) \\
Encoder & \texttt{allenai/scibert\_scivocab\_uncased} (768-dim) \\
Semantic Index & FAISS \texttt{IndexFlatIP} \\
Lexical Retriever & \texttt{BM25Okapi} (\texttt{rank-bm25}) \\
Fusion & Reciprocal Rank Fusion ($k=60$) \\
LLM Generator & \texttt{microsoft/Phi-3.5-mini-instruct} (3.8B param) \\
Max tokens gen. & 300 \\
Temperature & 0.3 (quasi deterministico) \\
\bottomrule
\end{tabular}
\end{center}

\textbf{Prompt template:}
\begin{codebox}
\begin{lstlisting}[style=python]
system_prompt = (
    "You are an expert NLP research assistant. "
    "Answer ONLY based on the provided context. "
    "Cite the papers you use."
)
context_text = "\n\n".join([doc[:200] for doc in retrieved_docs[:5]])
context_text = context_text[:3000]  # troncamento a 3000 char

messages = [
    {"role": "system", "content": system_prompt},
    {"role": "user",   "content": f"Context:\n{context_text}\n\nQuestion: {query}"}
]
\end{lstlisting}
\end{codebox}

%% ============================================================
\section{Metriche di Valutazione (Lab 5)}
%% ============================================================

\subsection{Precision@K}

\begin{formulabox}
\[
P@K = \frac{|\text{documenti rilevanti nei top-}K|}{K}
\]
\end{formulabox}
Un documento \`e ``rilevante'' se il titolo contiene almeno una keyword del Gold Standard.

\textbf{Gold Standard Lab 5} -- 5 query NLP con keyword attese:
\begin{itemize}
  \item ``BERT / self-attention'' $\to$ keyword: bert, attention, transformer, self-attention
  \item ``word2vec'' $\to$ keyword: word2vec, word embeddings, word vectors
  \item ``machine translation'' $\to$ keyword: translation, multilingual, parallel corpus
  \item ``topic modeling'' $\to$ keyword: topic, lda, latent dirichlet
  \item ``named entity recognition'' $\to$ keyword: ner, named entity, entity recognition
\end{itemize}

\subsection{Context Relevance (CR)}

\begin{formulabox}
\[
\text{CR} = \frac{1}{|R|} \sum_{d \in R} \cos\!\left(\text{emb}(q),\, \text{emb}(d[:200])\right)
\]
\end{formulabox}
Allineamento semantico medio tra query e snippet dei documenti recuperati.

\subsection{Answer Relevance (AR)}

\begin{formulabox}
\[
\text{AR} = \cos\!\left(\text{emb}(q),\, \text{emb}(\text{risposta})\right)
\]
\end{formulabox}
Misura se la risposta generata rimane centrata sulla query.

\subsection{Faithfulness (FF)}

\begin{formulabox}
\[
\text{FF} = \frac{|\{w \in \text{answer}: w \in \text{context}\}|}{|\text{answer\_content\_words}|}
\]
\end{formulabox}
Proxy lessicale: quante parole di contenuto della risposta compaiono nel contesto recuperato. Alta FF = risposta ancorata ai documenti (meno allucinazioni).

\subsection{Robustness}

\begin{formulabox}
\[
\text{Rob} = \frac{1}{|P|} \sum_{p \in P} \frac{|R_5(q) \cap R_5(p)|}{5}
\]
\end{formulabox}
Overlap tra risultati per query originale e parafrasi. Alta robustezza = retrieval stabile a riformulazioni.

%% ============================================================
\section{Risultati Lab 5: Esperimento Multi-Corpus}
%% ============================================================

\subsection{Disegno Sperimentale}

Variare la dimensione del corpus: $N \in \{1.000, 5.000, 15.000, 25.000, 44.949\}$

\subsection{Tabella Risultati Completa}

\begin{resultbox}
\begin{center}
\begin{tabular}{lccccc}
\toprule
\textbf{N} & \textbf{Metrica} & \textbf{Base (FAISS)} & \textbf{Hybrid (RRF)} & $\Delta$ & \textbf{\% Migl.}\\
\midrule
\multirow{4}{*}{1K}   & P@5   & 0.160 & 0.400 & +0.240 & +150\% \\
                       & CR    & 0.716 & 0.718 & +0.001 & trascur. \\
                       & AR    & 0.697 & 0.684 & -0.014 & -2\% \\
                       & FF    & 0.516 & 0.588 & +0.072 & +14\% \\
\midrule
\multirow{4}{*}{5K}   & P@5   & 0.240 & 0.560 & +0.320 & +133\% \\
                       & CR    & 0.724 & 0.714 & -0.010 & -1.4\% \\
                       & AR    & 0.687 & 0.674 & -0.014 & -2\% \\
                       & FF    & 0.572 & 0.667 & +0.095 & +16.6\% \\
\midrule
\multirow{4}{*}{15K}  & P@5   & 0.160 & 0.520 & +0.360 & +225\% \\
                       & CR    & 0.713 & 0.719 & +0.007 & +1\% \\
                       & AR    & 0.686 & 0.667 & -0.018 & -2.6\% \\
                       & FF    & 0.562 & 0.647 & +0.084 & +15\% \\
\midrule
\multirow{4}{*}{25K}  & P@5   & 0.200 & 0.440 & +0.240 & +120\% \\
                       & CR    & 0.720 & 0.730 & +0.010 & +1.4\% \\
                       & AR    & 0.679 & 0.682 & +0.003 & +0.4\% \\
                       & FF    & 0.560 & 0.587 & +0.027 & +4.8\% \\
\midrule
\multirow{4}{*}{44.9K} & P@5  & 0.120 & 0.360 & +0.240 & +200\% \\
                        & CR   & 0.719 & 0.726 & +0.007 & +1\% \\
                        & AR   & 0.678 & 0.666 & -0.012 & -1.8\% \\
                        & FF   & 0.496 & 0.569 & +0.074 & +14.9\% \\
\bottomrule
\end{tabular}
\end{center}
\end{resultbox}

\subsection{Analisi Critica dei Risultati}

\textbf{1. P@5: miglioramento sistematico del Hybrid}

L'approccio ibrido migliora consistentemente la P@5 su tutte le dimensioni corpus (+120\% a +225\%). BM25 recupera paper con acronimi tecnici esatti (NER, LDA, word2vec) che i soli embedding tendono a generalizzare.

\textbf{2. P@5 non monotona con la dimensione del corpus}
\begin{itemize}
  \item Picco a $N=5.000$ (P@5 hybrid = 0.560)
  \item Declino a $N=44.949$ (P@5 hybrid = 0.360)
  \item Spiegazione: corpus piccolo = paper rilevanti concentrati, alta densit\`a. Corpus grande = spazio raro, segnale di rilevanza diluito, pi\`u candidati competono
\end{itemize}

\textbf{3. CR e AR: differenze minime}

CR e AR sono costruiti su embedding SciBERT (usato da entrambi i sistemi). BM25 porta documenti lessicalmente precisi ma semanticamente pi\`u dispersi $\to$ leggero calo AR nel hybrid.

\textbf{4. Faithfulness: miglioramento consistente}

I documenti BM25 contengono i termini esatti della query $\to$ il LLM li riusa nell'answer $\to$ maggiore overlap lessicale answer-context. Questo \`e particolarmente utile nei domini tecnici dove il vocabolario \`e vincolato.

\textbf{5. Robustness: declino catastrofico con la scala}

\begin{center}
\begin{tabular}{lc}
\toprule
\textbf{N docs} & \textbf{Robustness (overlap@5)} \\
\midrule
1.000   & 0.233 (23.3\%) \\
5.000   & 0.100 (10.0\%) \\
15.000  & 0.033 (3.3\%) \\
25.000  & 0.000 (0.0\%) \\
44.949  & 0.033 (3.3\%) \\
\bottomrule
\end{tabular}
\end{center}

In corpus piccoli pochi paper dominanti $\to$ le parafrasi recuperano gli stessi. In corpus grandi lo spazio \`e denso di candidati vicini $\to$ una minima riformulazione cambia completamente i top-5. Soluzioni: query expansion, ensemble retriever, re-ranking gerarchico.

\textbf{6. Precision per query (hybrid)}
\begin{center}
\begin{tabular}{lccccc}
\toprule
\textbf{Query} & 1K & 5K & 15K & 25K & 44.9K \\
\midrule
BERT/self-attention & 0.20 & 0.40 & 0.40 & 0.00 & 0.20 \\
word2vec            & 0.40 & 0.60 & 0.60 & 0.60 & 0.40 \\
machine translation & 0.20 & 0.40 & 0.60 & 0.60 & 0.40 \\
topic modeling      & 0.40 & 0.60 & 0.20 & 0.20 & 0.20 \\
named entity recog. & 0.80 & 0.80 & 0.80 & 0.80 & 0.60 \\
\midrule
\textbf{Media}      & \textbf{0.40} & \textbf{0.56} & \textbf{0.52} & \textbf{0.44} & \textbf{0.36} \\
\bottomrule
\end{tabular}
\end{center}

``Named entity recognition'' \`e il top performer stabile: acronimo NER molto specifico, BM25 lo identifica bene. ``BERT/self-attention'' instabile (0.00 a 25K): entrambi i termini sono ubiqui nel corpus NLP.

%% ============================================================
\section{Confronto Approcci: Tabella Sinottica}
%% ============================================================

\begin{center}
\begin{tabular}{p{2.8cm}p{3.5cm}p{3.5cm}p{3.5cm}}
\toprule
\textbf{Aspetto} & \textbf{TF-IDF/Jaccard (Lab 1-2)} & \textbf{BERTopic (Lab 4)} & \textbf{RAG Hybrid (Lab 5)} \\
\midrule
\textbf{Task} & Misura similarit\`a definizioni; recupero onomasiologico & Scoperta automatica topic & Risposta a domande su corpus \\
\midrule
\textbf{Rappresentazione} & Vettori sparsi (BoW / TF-IDF) & Embedding densi + riduzione dim & Embedding densi (FAISS) + sparsi (BM25) \\
\midrule
\textbf{Semantica} & Limitata (no sinonimi) & Alta (SentenceTransformer) & Alta semantica + alta precisione lessicale \\
\midrule
\textbf{Scalabilit\`a} & Alta, veloce & Media (UMAP/HDBSCAN) & Dipende da corpus (best: 5K-15K) \\
\midrule
\textbf{Output} & Score similarit\`a / synset & Lista topic + keywords & Risposta in linguaggio naturale \\
\midrule
\textbf{Limite principale} & Manca comprensione semantica profonda & Config errata $\to$ failure mode & Robustness degrada con scala \\
\bottomrule
\end{tabular}
\end{center}

%% ============================================================
\section{Teoria: Embedding e Retrieval}
%% ============================================================

\subsection{Dense vs. Sparse Representation}

\begin{defbox}
\textbf{Rappresentazioni sparse} (TF-IDF, BM25):
\begin{itemize}
  \item Dimensione = vocabolario (tipicamente $10^4$--$10^6$)
  \item Quasi tutte le componenti sono zero
  \item Interpretabili: ogni dimensione = un termine
  \item Veloci, nessuna inferenza neurale
\end{itemize}

\textbf{Rappresentazioni dense} (BERT, SentenceTransformer):
\begin{itemize}
  \item Dimensione fissa ridotta (128--1024)
  \item Tutte le componenti non-zero
  \item Catturano significato distribuito, sinonimi, relazioni implicite
  \item Pi\`u lente (require inferenza del modello)
\end{itemize}
\end{defbox}

\subsection{Sentence Transformers}

\begin{defbox}
\textbf{SentenceTransformer}: modello BERT fine-tuned per produrre rappresentazioni a livello di frase.
\begin{itemize}
  \item Training Siamese/triplet: frasi simili vicine, dissimili lontane nello spazio vettoriale
  \item Pooling sull'ultimo layer per ottenere vettore singolo per la frase
  \item Modelli comuni: \texttt{all-MiniLM-L6-v2} (384-dim, leggero), \texttt{all-mpnet-base-v2} (768-dim, alta qualit\`a)
\end{itemize}
\end{defbox}

\subsection{SciBERT vs. BERT generale}

SciBERT (\texttt{allenai/scibert\_scivocab\_uncased}):
\begin{itemize}
  \item Stessa architettura BERT (12 layer, 768-dim base)
  \item Vocabolario costruito su testi scientifici (formule, termini tecnici, notazione matematica)
  \item Pre-training su 1.14M paper da Semantic Scholar
  \item Superiore su IR scientifico rispetto a BERT generale
\end{itemize}

\subsection{Cos'\`e HDBSCAN: Intuizione}

HDBSCAN costruisce una gerarchia di cluster basata sulla densit\`a:
\begin{enumerate}
  \item Calcola la core distance di ogni punto (distanza al $k$-esimo vicino)
  \item Costruisce il minimum spanning tree con distanze di mutua raggiungibilit\`a
  \item Condensa la gerarchia: rimuove rami che ``si dividono'' troppo rapidamente
  \item Estrae cluster stabili con Excess of Mass (\texttt{eom})
  \item Punti non assegnati a nessun cluster = outlier ($-1$)
\end{enumerate}

Vantaggio fondamentale per BERTopic: non si specifica quanti topic esistono, emergono dai dati.

%% ============================================================
\section{Riepilogo Formule}
%% ============================================================

\begin{center}
\renewcommand{\arraystretch}{1.5}
\begin{tabular}{lll}
\toprule
\textbf{Formula} & \textbf{Espressione} & \textbf{Uso} \\
\midrule
Jaccard (SimLex) & $|A \cap B| / |A \cup B|$ & Overlap lessicale definizioni \\
Cosine Similarity & $\vec{a} \cdot \vec{b} \,/\, (||\vec{a}|| \cdot ||\vec{b}||)$ & SimSem, FAISS, CR, AR \\
TF-IDF & $\text{tf}(t,d) \times \log(N/\text{df}(t))$ & Vettorizzazione Lab1-2 \\
BM25 score & $\sum_i \text{IDF}(q_i) \cdot \frac{f(q_i,D)(k_1+1)}{f(q_i,D)+k_1(1-b+b|D|/\text{avg})}$ & Retrieval lessicale \\
RRF score & $\sum_{r} 1/(k + \text{rank}_r(d))$ & Fusione ranking ibrido \\
c-TF-IDF & $\frac{f(t,c)}{|c|} \cdot \log\frac{N}{\text{df}(t)}$ & Topic keywords (BERTopic) \\
Precision@K & $|\text{rilevanti} \cap \text{top-}K| \,/\, K$ & Qualit\`a retrieval \\
Faithfulness & $|\{w \in \text{ans}: w \in \text{ctx}\}| / |\text{ans}|$ & Grounding LLM \\
Robustness & $\frac{1}{|P|}\sum_p |R_5(q) \cap R_5(p)| \,/\, 5$ & Stabilit\`a a parafrasi \\
Context Rel. & $\frac{1}{|R|}\sum_d \cos(q, d[:200])$ & Pertinenza contesto \\
Answer Rel. & $\cos(\text{emb}(q), \text{emb}(\text{ans}))$ & Centramento risposta \\
\bottomrule
\end{tabular}
\end{center}

%% ============================================================
\section{Domande Tipiche da Esame}
%% ============================================================

\begin{enumerate}
  \item \textbf{Differenza tra SimLex e SimSem}: Jaccard misura overlap lessicale puro su insiemi di parole; TF-IDF coseno misura prossimit\`a semantica in spazio vettoriale. Per concetti astratti SimSem > SimLex perch\'e cattura sinonimi.

  \item \textbf{Cos'\`e la ricerca onomasiologica?} Data una definizione, recuperare il synset WordNet corretto. Inverso della consultazione dizionario. Implementata con TF-IDF coseno tra definizione e glosse candidate.

  \item \textbf{Perch\'e i concetti concreti sono pi\`u recuperabili?} Il referente fisico vincola il vocabolario: persone diverse usano parole simili per descriverlo. I concetti astratti hanno alta variabilit\`a e dipendenza culturale.

  \item \textbf{Descrivere la pipeline BERTopic}: Documenti $\to$ Embedding (SentenceTransformer) $\to$ Riduzione dimensionale (UMAP) $\to$ Clustering (HDBSCAN) $\to$ Estrazione keyword (c-TF-IDF).

  \item \textbf{Perch\'e la Config 3 di BERTopic fallisce?} \texttt{min\_cluster\_size=150} crea cluster troppo grandi e tematicamente eterogenei. Il c-TF-IDF non riesce a estrarre segnale discriminativo e le stopwords emergono nel top-5.

  \item \textbf{Cos'\`e RRF e perch\'e si usa?} Reciprocal Rank Fusion combina ranking da retriever diversi (FAISS + BM25) senza richiedere normalizzazione degli score: $\sum_r 1/(k+\text{rank}_r(d))$. Sfrutta la complementariet\`a di retrieval semantico e lessicale.

  \item \textbf{Trend P@5 vs. corpus size}: non monotono, picco a 5K. Corpus troppo piccolo = pochi documenti rilevanti. Corpus troppo grande = segnale di rilevanza diluito in spazio denso di candidati.

  \item \textbf{Perch\'e la robustezza crolla a scala?} In corpus grandi lo spazio embedding \`e molto denso: una minima riformulazione della query sposta i top-k verso documenti simili ma diversi. Fix: query expansion, ensemble retrieval, re-ranking.

  \item \textbf{Differenza BM25 vs. TF-IDF}: BM25 aggiunge saturazione della frequenza dei termini ($k_1$) e normalizzazione per lunghezza del documento ($b$), migliorando il ranking rispetto al semplice TF-IDF.

  \item \textbf{Cosa misura la Faithfulness?} La percentuale di parole di contenuto della risposta LLM che compaiono nel contesto recuperato. Proxy lessicale per l'assenza di allucinazioni.
\end{enumerate}

\end{document}
